% Sections 13.1 to 13.6, from lectures/L13/appendix/a_theory.md.

\section{Derivata för trigonometriska funktioner}
\label{c13:sec:trig}

\begin{narrowtable}{@{}ll@{}}
  \thead{Funktion} & \thead{Derivata} \\
  \midrule
  $\sin(x)$ & $\cos(x)$ \\
  $\cos(x)$ & $-\sin(x)$ \\
  $\sin(kx)$ & $k\cos(kx)$ \\
  $\cos(kx)$ & $-k\sin(kx)$ \\
\end{narrowtable}

\begin{exempel}
\[
  f(x) = 3\sin(2x) \quad \Rightarrow \quad f'(x) = 3 \cdot 2\cos(2x) = 6\cos(2x)
\]
\end{exempel}

\section{Derivata för exponentialfunktioner}
\label{c13:sec:exp}

\begin{narrowtable}{@{}ll@{}}
  \thead{Funktion} & \thead{Derivata} \\
  \midrule
  $e^x$ & $e^x$ \\
  $e^{kx}$ & $ke^{kx}$ \\
  $a^x$ & $a^x \ln a$ \\
\end{narrowtable}

\textbf{Härledning av $\frac{d}{dx}e^x = e^x$.} Med derivatans definition från
\secref{c12:sec:definition} blir
\[
  \lim_{h \to 0} \frac{e^{x+h} - e^x}{h} = e^x \lim_{h \to 0} \frac{e^h - 1}{h} = e^x \cdot 1
  = e^x.
\]

\section{Derivata för logaritmfunktioner}
\label{c13:sec:log}

{\renewcommand{\arraystretch}{2}%
\begin{narrowtable}{@{}ll@{}}
  \thead{Funktion} & \thead{Derivata} \\
  \midrule
  $\ln(x)$ & $\dfrac{1}{x}$ \\
  $\ln(kx)$ & $\dfrac{1}{x}$ \\
  $\log_{10}(x)$ & $\dfrac{1}{x \ln 10}$ \\
\end{narrowtable}}

\begin{aside}
\note[Obs] $\ln(kx) = \ln k + \ln x$, vars derivata är $1/x$: konstanten $\ln k$ försvinner.
\end{aside}

\section{Produktregeln och kedjeregeln}
\label{c13:sec:regler}
Tabellerna ovan ger derivatan av en enskild funktion. När två funktioner multipliceras med
varandra, eller när en funktion sätts in i en annan, behövs två regler till.

\subsection{Produktregeln}
\label{c13:sec:produktregeln}
Derivatan av en produkt $f(x) = g(x) \cdot h(x)$ är
\[
  f'(x) = g'(x) \cdot h(x) + g(x) \cdot h'(x).
\]
Derivera alltså en faktor i taget medan den andra står kvar, och addera de två termerna.

\begin{exempel}
$f(x) = x^2 \sin(x)$ med $g(x) = x^2$ och $h(x) = \sin(x)$ ger
\[
  f'(x) = 2x \cdot \sin(x) + x^2 \cdot \cos(x).
\]
\end{exempel}

\subsection{Kedjeregeln}
\label{c13:sec:kedjeregeln}
En sammansatt funktion $f(x) = g(u(x))$ består av en \textbf{yttre funktion} $g$ och en
\textbf{inre funktion} $u(x)$. Derivatan är den yttre funktionens derivata, beräknad i den inre
funktionen, gånger den inre funktionens derivata:
\[
  f'(x) = g'(u(x)) \cdot u'(x).
\]

\begin{exempel}
$f(x) = \ln(2x + 1)$ har den yttre funktionen $\ln u$ och den inre funktionen $u = 2x + 1$, med
$u' = 2$:
\[
  f'(x) = \frac{1}{2x + 1} \cdot 2 = \frac{2}{2x + 1}.
\]
\end{exempel}

Tabellraderna för $\sin(kx)$, $\cos(kx)$, $e^{kx}$ och $\ln(kx)$ ovan är specialfall av
kedjeregeln, med den inre funktionen $u = kx$ och $u' = k$. Till exempel är
\[
  \frac{d}{dx}\sin(kx) = \cos(kx) \cdot k = k\cos(kx), \qquad
  \frac{d}{dx}\ln(kx) = \frac{1}{kx} \cdot k = \frac{1}{x}.
\]

\section{Samlad derivatatabell}
\label{c13:sec:tabell}

\begin{narrowtable}{@{}ll@{}}
  \thead{Funktion $f(x)$} & \thead{Derivata $f'(x)$} \\
  \midrule
  $x^n$ & $nx^{n-1}$ \\
  $e^x$ & $e^x$ \\
  $e^{kx}$ & $ke^{kx}$ \\
  $\ln x$ & $1/x$ \\
  $\sin x$ & $\cos x$ \\
  $\cos x$ & $-\sin x$ \\
  $\sin(kx)$ & $k\cos(kx)$ \\
  $\cos(kx)$ & $-k\sin(kx)$ \\
\end{narrowtable}

\section{Typexempel}
\label{c13:sec:typexempel}

\begin{exempel}[Derivata av vanliga funktioner]
Derivera \textbf{a)} $f(x) = 5e^{3x}$, \textbf{b)} $f(x) = \ln(4x)$ och
\textbf{c)} $f(x) = 2\cos(x) - 3\sin(x)$.
\begin{align*}
  \textbf{a)}\quad f'(x) &= 5 \cdot 3e^{3x} = 15e^{3x} \\
  \textbf{b)}\quad f'(x) &= 1/x \\
  \textbf{c)}\quad f'(x) &= -2\sin(x) - 3\cos(x)
\end{align*}
\end{exempel}

\begin{exempel}[Beräkna derivata i en punkt]
Strömmen $i(t) = I_0 e^{-t/\tau}$ ampere beskriver urladdningen av en kondensator. Beräkna
$i'(0)$.
\[
  i'(t) = -\frac{I_0}{\tau} e^{-t/\tau} \quad \Rightarrow \quad i'(0) = -\frac{I_0}{\tau}
\]
$i'(0)$ är den momentana strömförändringshastigheten vid $t = 0$.
\end{exempel}

\begin{exempel}[Stationär punkt för exponentialfunktion]
Bestäm extrempunkten för $f(x) = xe^{-x}$. Produktregeln (\secref{c13:sec:produktregeln})
med $g(x) = x$ och $h(x) = e^{-x}$ ger
\[
  f'(x) = e^{-x} + x \cdot (-e^{-x}) = e^{-x}(1 - x)
\]
$f'(x) = 0$ ger $x = 1$, eftersom $e^{-x} \neq 0$. Andraderivatan är
\[
  f''(x) = -e^{-x}(1 - x) + e^{-x}(-1) = e^{-x}(x - 2),
\]
och
\[
  f''(1) = e^{-1}(1 - 2) = -e^{-1} < 0 \quad \Rightarrow \quad
  \text{\textbf{maximum} i } (1, e^{-1}) \approx (1;\, 0{,}37).
\]
\end{exempel}
